\section{相关工作}
\label{sec:related_work}

\subsection{面向机器人策略的流匹配}

流匹配~\cite{lipman2023flow,liu2023flow}学习一个时间相关的速度场$v_\theta(x_t, t)$，通过ODE $\mathrm{d}x_t = v_\theta(x_t, t)\,\mathrm{d}t$将样本从先验分布传输到数据分布。与扩散模型~\cite{ho2020denoising,song2021score}相比，流匹配具有更简单的训练过程（无需噪声调度）、确定性采样，以及与最优传输（OT-CFM）~\cite{tong2024improving}结合时更直的传输路径。

在机器人领域，Diffusion Policy~\cite{chi2023diffusion}首次证明去噪扩散可以有效建模视觉运动控制的多模态动作分布，引入了动作分块（action chunking）范式。VITA~\cite{gao2026vita}将这一思想扩展到流匹配，提出从视觉观测隐变量直接流向动作隐变量——消除了随机噪声起始点，实现了自然的多模态动作生成。$\pi_0$~\cite{black2024pi0}将流匹配策略扩展到基础模型级别，通过视觉-语言-动作预训练展示了跨任务的强零样本泛化能力。然而，所有这些方法在视觉编码阶段都采用全局平均池化，丢失了精密任务所需的空间结构信息。本文的工作建立在VITA架构基础之上，通过自监督空间注意力对其进行增强。

\subsection{观测堆叠与帧差分}

最简单的时间编码策略是\emph{观测堆叠}：将最近$k$帧拼接作为输入~\cite{mnih2015human}。虽然直接简单，但这仅提供局部时间上下文，未显式编码帧间\emph{变化内容}，且随时间窗长度增加扩展性差。我们的特征差分方法（第~\ref{sec:temporal_diff}节）通过计算显式特征空间差分来解决这一问题，在不随观测窗增加维度的情况下编码"语义速度"。这种差分表示也为\fasa{}的变化检测提供了自然的基础。

\subsection{视觉运动策略中的空间注意力}

空间注意力机制旨在引导视觉编码器关注任务相关区域，而非平等对待所有空间位置。

\paragraph{学习的空间表示。}Spatial Softmax~\cite{levine2016end}从卷积特征图计算逐通道的期望坐标，提供紧凑的空间表示，已成为视觉运动策略学习的标准工具。然而，它对所有通道施加均匀注意力，不区分任务相关和无关区域。

\paragraph{通道与空间注意力模块。}Squeeze-and-Excitation（SE）网络~\cite{hu2018squeeze}通过全局池化和轻量级MLP学习通道注意力权重，按重要性重新加权特征通道。CBAM~\cite{woo2018cbam}通过额外的空间注意力分支学习二维注意力图进行扩展。Coordinate Attention~\cite{hou2021coordinate}将通道注意力分解为两个空间方向，捕获沿各轴的长程位置依赖关系。虽然这些模块对一般视觉任务有效，但它们从数据分布中学习\emph{静态}注意力模式，不会适应episode内当前的交互动态。

\paragraph{基于基础模型的知识蒸馏。}视觉基础模型如Grounding DINO~\cite{liu2023grounding}和CLIP~\cite{radford2021learning}编码了丰富的空间和语义知识，可通过注意力热图匹配或特征蒸馏损失传递给轻量级主干网络。虽然有效，但这些方法引入了对特定教师模型的依赖，需要离线特征提取流水线，且产生不随任务阶段变化的静态空间目标。

我们的\fasa{}模块采取了根本不同的方法：它不依赖外部空间监督或学习静态注意力模式，而是通过帧间特征时间对比\emph{自监督}地发现任务相关区域。核心洞察是，连续帧之间视觉变化最大的区域对应于活跃的交互区——恰恰是操作任务中最相关的区域。这使得\fasa{}的注意力天然具有\emph{动态}性，随着任务从接近阶段到跟踪阶段实时适应。
